\documentclass{article}
\title{A Tiny Sample Paper for Anatomy Decomposition Tests}
\author{Test Author}
\begin{document}
\maketitle

\section{Introduction}

This is the introduction. It cites \cite{foo23} and \cite{bar22}
as motivation.

A second paragraph in the introduction. It cites \cite{baz21}.

\section{Background}

Background body text.

\subsection{Bitemporal Storage}

A subsection about bitemporal storage.
We rely on the prior work of \cite{snodgrass95}.

\begin{equation}
\label{eqn:basic}
y = mx + b
\end{equation}

\begin{theorem}
\label{thm:halting}
For every Turing machine $M$ and input $w$, the question of
whether $M$ halts on $w$ is undecidable.
\end{theorem}

\begin{proof}
By diagonalization. Suppose for contradiction that a decider $H$
exists. See \cite{turing36}.
\end{proof}

\section{Method}

\begin{figure}
\label{fig:arch}
\caption{Architecture diagram of the system.}
\end{figure}

\begin{table}
\label{tbl:perf}
\caption{Performance comparison across benchmarks.}
\end{table}

\begin{algorithm}
\label{alg:probe}
\caption{Probe routine.}
\textbf{Input:} dataset $D$, threshold $\tau$.
\textbf{Output:} accept / reject.
For each $x \in D$, evaluate $f(x)$; return accept iff
$\sum f(x) \geq \tau$.
\end{algorithm}

\section{Conclusion}

We have shown that anatomy decomposition works. Future work in
\cite{future26}.

\end{document}
